@@M0@@ \uline{is exact).}

Put @@M1@@, @@M2@@, and define @@M3@@ recursively. We verify at once that @@M4@@. It remains to prove that, for @@M5@@ sufficiently large, @@M6@@. Indeed, consider the decreasing sequence of @@M7@@-modules @@M8@@, and put @@M9@@. The decreasing sequence of closed subsets @@M10@@ of @@M11@@ is stationary; let @@M12@@ be the stationary value. Let us show that @@M13@@. Otherwise, let @@M14@@ be a maximal point of @@M15@@. After replacing @@M16@@ by @@M17@@, we may suppose that @@M18@@ is local with closed point @@M19@@. By the choice of @@M20@@, @@M21@@ is artinian for @@M22@@ sufficiently large, so the sequence @@M23@@ is stationary. Thus there exists @@M24@@ such that @@M25@@. Consequently, @@M26@@ is a monomorphism, hence a closed immersion (Nakayama); since @@M27@@, we have @@M28@@ and @@M29@@, a contradiction.

\sgaStableTarget{sga6.e12.sec.3}\section*{3. Proof of I: Oort's dévissage}

\sgaStableTarget{sga6.e12.prop.3-1}\textbf{Proposition 3.1.} (Oort, Bull. soc. math. France, vol. 90-1962). \uline{Let} @@M30@@ \uline{be a locally noetherian scheme,} @@M31@@ \uline{a proper morphism, and} @@M32@@ \uline{a closed subscheme. Suppose that:}
\begin{enumerate}[label=\roman*)]
\item @@M33@@ \uline{is defined by an ideal} @@M34@@ \uline{contained in the nilradical} @@M35@@ \uline{of} @@M36@@ \uline{and} @@M37@@.
\item @@M38@@ \uline{is flat over} @@M39@@ \uline{and cohomologically flat over} @@M40@@ \uline{in dimension} @@M41@@ (EGA III 7.8.1).
\item @@M42@@ \uline{is flat over} @@M43@@ \uline{and cohomologically flat over} @@M44@@ \uline{in dimensions} @@M45@@ \uline{and} @@M46@@\uline{, so that (by} EGA III 7.4.6\uline{) there are a coherent} @@M47@@\uline{-module} @@M48@@ \uline{and an isomorphism, functorial in the coherent} @@M49@@\uline{-module} @@M50@@,

